\chapter{Random Groups and the Cohen--Lenstra Heuristics}\label{ch:cohenlenstra}
\begin{center}\itshape
A group with many symmetries is a group that nature is shy to produce.
\end{center}

Gauss computed class numbers by the hundred in the \emph{Disquisitiones Arithmeticae} (1801), and the tables of class groups of imaginary quadratic fields have grown ever since. Look at the part of the class group of order a power of $3$: for about $56\%$ of the discriminants it is trivial, for about $28\%$ it is cyclic of order $3$, and the group $\Z/3\times\Z/3$ is much rarer than $\Z/9$, although both have $9$ elements. In 1983 Henri Cohen and Hendrik Lenstra proposed an explanation at the Journées Arithmétiques of Noordwijkerhout. Their principle: a finite abelian $p$-group $G$ should occur with probability proportional to $\frac1{|\operatorname{Aut}G|}$. The group $\Z/3\times\Z/3$ has $48$ automorphisms, $\Z/9$ only $6$, and the ratio of the observed frequencies is close to $8$.

The heuristic was a guess, weighted by symmetry, and it predicted everything in the tables. It also had a mechanism hidden inside it, found by Eduardo Friedman and Lawrence Washington in 1989: the Cohen--Lenstra law is the law of the cokernel of a large random matrix with entries in the $p$-adic integers. For function fields this is no metaphor, since the class group of a curve over $\mathbb F_q$ is the cokernel of $1-F$, with $F$ the Frobenius acting on a lattice.

This chapter proves the probabilistic side of the story completely. We compute the moments $\mathbb E\#\operatorname{Sur}(X,G)$ of random cokernels, all equal to $1$ in the limit; we derive the Cohen--Lenstra law from them, and we prove the theorem of Melanie Wood that moments of this size determine the law of a random group. Her theorem, the moment method for random groups, has become the main tool of the subject: to prove that some arithmetic groups are Cohen--Lenstra distributed, it suffices to compute their moments. The next chapter computes one such moment for class groups: the average number of $3$-torsion elements, by Davenport and Heilbronn.

\begin{reminder}[from Tomes~2, 3 and~12]
\emph{Finite abelian groups} (Tome~2, Theorem~7.6). A finite abelian $p$-group is isomorphic to $\bigoplus_i\Z/p^{a_i}$ for a unique partition $a_1\ge a_2\ge\cdots$, its \emph{type}.
\emph{Class groups} (Tome~12, Chapter~2). The class group $\operatorname{Cl}(K)$ of a number field is a finite abelian group; for an imaginary quadratic field $\Q(\sqrt{-d})$ its order is the class number $h(-d)$, the number of reduced binary quadratic forms of discriminant $-d$.
\emph{Convergence in law} (Tome~3). For random variables with values in a countable set, convergence in law means $\mathbb P(X_n=x)\to\mathbb P(X=x)$ for each $x$; if the family is tight, every subsequence has a further subsequence converging in law.
\end{reminder}

Throughout, $p$ is a fixed prime, all groups are finite abelian $p$-groups, written additively, and $\operatorname{Sur}(A,G)$, $\operatorname{Hom}(A,G)$, $\operatorname{Aut}(G)$ have their usual meanings. The constant
\[
c_\infty=\prod_{i\ge1}\bigl(1-p^{-i}\bigr)
\]
appears everywhere; $c_\infty=0.2888$ for $p=2$, $0.5601$ for $p=3$, $0.7603$ for $p=5$.

\section{Counting homomorphisms}

For a group $A$ of type $a_1\ge a_2\ge\cdots$, let $r_k(A)=\#\{i:a_i\ge k\}$, the \emph{conjugate type}. Equivalently $r_k(A)=\dim_{\mathbb F_p}p^{k-1}A/p^kA$, and $r_1(A)$, the number of cyclic factors, is the \emph{$p$-rank} of $A$. A group of exponent dividing $p^e$ is determined by the vector $(r_1,\dots,r_e)$, which can be any non-increasing vector of non-negative integers.

\begin{lemma}\label{lem:homcount}
For groups $A$, $G$, $|\operatorname{Hom}(A,G)|=p^{\sum_kr_k(A)r_k(G)}$ and $|\wedge^2G|=p^{\sum_k\binom{r_k(G)}2}$. In particular $|\operatorname{End}(G)|=p^{\sum_kr_k(G)^2}$.
\end{lemma}
\begin{proof}
$\operatorname{Hom}(\Z/p^a,\Z/p^b)\cong\Z/p^{\min(a,b)}$, so $|\operatorname{Hom}(A,G)|=p^{\sum_{i,j}\min(a_i,g_j)}$, and $\min(a,g)=\#\{k\ge1:a\ge k\text{ and }g\ge k\}$ turns the exponent into $\sum_kr_k(A)r_k(G)$. Similarly $\wedge^2(\bigoplus_j\Z/p^{g_j})\cong\bigoplus_{j<l}\Z/p^{\min(g_j,g_l)}$, of order $p^{\sum_{j<l}\min(g_j,g_l)}$, and $\sum_{j<l}\min(g_j,g_l)=\sum_k\binom{r_k(G)}2$.
\end{proof}

\begin{lemma}\label{lem:autcount}
Let $G$ have type with $m_a$ parts equal to $a$, for each $a$. Then
\[
|\operatorname{Aut}G|=p^{\sum_kr_k(G)^2}\prod_a\prod_{i=1}^{m_a}\bigl(1-p^{-i}\bigr)\ \ge\ c_\infty^{\,e}\,p^{\sum_kr_k(G)^2}
\]
if $G$ has exponent dividing $p^e$.
\end{lemma}
\begin{proof}
An endomorphism $\varphi$ of $G$ is an automorphism if and only if the induced map $\bar\varphi$ on $G/pG$ is bijective. Indeed, if $\bar\varphi$ is onto, the image $I$ of $\varphi$ satisfies $I+pG=G$, so the $p$-group $G/I$ equals $p(G/I)$ and is trivial; then $\varphi$ is onto, hence bijective as $G$ is finite. The converse is clear.

Write $G=\bigoplus_j\Z/p^{g_j}$ with generators $\epsilon_j$, and $\varphi(\epsilon_j)=\sum_lc_{lj}\epsilon_l$, where $c_{lj}$ is any element of $p^{\max(0,g_l-g_j)}\Z/p^{g_l}$, the condition for $\epsilon_j\mapsto c_{lj}\epsilon_l$ to be well defined. A uniform random endomorphism is given by independent uniform $c_{lj}$ in these sets. Modulo $p$, $\bar c_{lj}=0$ when $g_l>g_j$, and $\bar c_{lj}$ is uniform on $\mathbb F_p$ when $g_l=g_j$. Ordering the $j$ by $g_j$, the matrix of $\bar\varphi$ is block triangular, with independent uniform diagonal blocks of sizes $m_a$. It is invertible if and only if each diagonal block is, which has probability $\prod_a\frac{|GL_{m_a}(\mathbb F_p)|}{p^{m_a^2}}=\prod_a\prod_{i=1}^{m_a}(1-p^{-i})$. Multiply by $|\operatorname{End}G|$ (Lemma~\ref{lem:homcount}). There are at most $e$ distinct part sizes, and each product over $i$ is at least $c_\infty$.
\end{proof}

\begin{lemma}\label{lem:subgroupsum}
For $G$ of exponent dividing $p^e$,
\[
\sum_{H\le G}|\wedge^2H|\le\Bigl(\frac3{c_\infty}\Bigr)^e|\wedge^2G|.
\]
In particular $G$ has at most $(\frac3{c_\infty})^e|\wedge^2G|$ subgroups.
\end{lemma}
\begin{proof}
Write $\lambda_k=r_k(G)$. A subgroup $H$ of conjugate type $\nu=(\nu_1,\dots,\nu_e)$ is the image of an injective homomorphism from the abstract group $H_\nu$ of that type, and two injections have the same image exactly when they differ by an automorphism of $H_\nu$. Hence, by Lemmas~\ref{lem:homcount} and~\ref{lem:autcount}, the number of such $H$ is at most
\[
\frac{|\operatorname{Hom}(H_\nu,G)|}{|\operatorname{Aut}H_\nu|}\le c_\infty^{-e}p^{\sum_k(\nu_k\lambda_k-\nu_k^2)},
\]
and this vanishes unless $\nu_k\le\lambda_k$ for every $k$. Indeed, in a finite abelian $p$-group $A$ the number of cyclic factors equals $\dim_{\mathbb F_p}A[p]$, where $A[p]=\{a:pa=0\}$, since $(\Z/p^a)[p]$ has order $p$. Applying this to $A=p^{k-1}H$, whose cyclic factors are the $\Z/p^{h_i-k+1}$ with $h_i\ge k$, gives $\nu_k=\dim(p^{k-1}H)[p]$; likewise $\lambda_k=\dim(p^{k-1}G)[p]$, and $(p^{k-1}H)[p]\subseteq(p^{k-1}G)[p]$. Each such $H$ has $|\wedge^2H|=p^{\sum_k\binom{\nu_k}2}$, and
\[
\nu\lambda-\nu^2+\binom\nu2=\binom\lambda2-\binom{\lambda-\nu}2.
\]
Since $\sum_{t\ge0}p^{-\binom t2}\le1+1+\frac12+\frac18+\dots\le3$, summing over $0\le\nu_k\le\lambda_k$ independently for each $k$ gives the claim. The second statement follows from $|\wedge^2H|\ge1$.
\end{proof}

\section{Cokernels of random matrices}

Fix $e\ge1$ and let $R=\Z/p^e$. For an $n\times n$ matrix $M$ over $R$, its \emph{cokernel} is $\operatorname{cok}M=R^n/MR^n$, a group of exponent dividing $p^e$. A group $G$ of exponent dividing $p^e$ is an $R$-module, and homomorphisms of groups between such groups are $R$-linear.

\begin{proposition}\label{prop:freesur}
For $G$ of exponent dividing $p^e$ and $p$-rank $r\le n$,
\[
\#\operatorname{Sur}(R^n,G)=|G|^n\prod_{i=0}^{r-1}\bigl(1-p^{i-n}\bigr).
\]
\end{proposition}
\begin{proof}
A homomorphism $R^n\to G$ is the choice of the images $g_1,\dots,g_n$ of the basis vectors, and it is onto if and only if the $g_i$ generate $G$, if and only if their images generate $G/pG\cong\mathbb F_p^r$ (as in the proof of Lemma~\ref{lem:autcount}). The map $G^n\to(G/pG)^n$ is $|pG|^n$-to-one, and the number of $n$-tuples spanning $\mathbb F_p^r$ is the number of surjective linear maps $\mathbb F_p^n\to\mathbb F_p^r$, that is $(p^n-1)(p^n-p)\cdots(p^n-p^{r-1})$ (choose the images of the $r$ coordinate forms of a dual basis independently). Multiplying, $|pG|^np^{nr}\prod_{i<r}(1-p^{i-n})=|G|^n\prod_{i<r}(1-p^{i-n})$.
\end{proof}

\begin{theorem}[moments of random cokernels]\label{thm:cokmoments}
Let $M$ be a uniform random $n\times n$ matrix over $R=\Z/p^e$. For every $G$ of exponent dividing $p^e$ and $p$-rank $r\le n$,
\[
\mathbb E\,\#\operatorname{Sur}(\operatorname{cok}M,G)=\prod_{i=0}^{r-1}\bigl(1-p^{i-n}\bigr)\xrightarrow[n\to\infty]{}1.
\]
\end{theorem}
\begin{proof}
Surjections $\operatorname{cok}M\to G$ correspond to surjections $F:R^n\to G$ with $F\circ M=0$, that is with $F(Mx)=0$ for all $x$, that is $F(m_j)=0$ for each column $m_j$ of $M$. For a fixed surjective $F$, $F(m)$ is uniform on $G$ when $m$ is uniform on $R^n$ (each fibre of the onto homomorphism has the same size), and the $n$ columns are independent, so $\mathbb P(F\circ M=0)=|G|^{-n}$. Summing over $F$ and using Proposition~\ref{prop:freesur} gives the formula.
\end{proof}

So, whatever the group $G$, a large random cokernel has on average one surjection onto $G$. This is all the information we need, as the rest of the chapter shows. First, the law itself.

\begin{theorem}[Friedman--Washington, 1989]\label{thm:fw}
Let $M$ be a uniform random $n\times n$ matrix over $R=\Z/p^e$, and $G$ a group of exponent dividing $p^{e-1}$ and $p$-rank $r\le n$. Then
\[
\mathbb P(\operatorname{cok}M\cong G)=\frac1{|\operatorname{Aut}G|}\prod_{i=0}^{r-1}\bigl(1-p^{i-n}\bigr)\prod_{j=1}^n\bigl(1-p^{-j}\bigr)\xrightarrow[n\to\infty]{}\frac{c_\infty}{|\operatorname{Aut}G|}.
\]
\end{theorem}
\begin{proof}
The submodules $K\subseteq R^n$ with $R^n/K\cong G$ are the kernels of the surjections $R^n\to G$, and two surjections have the same kernel exactly when they differ by an automorphism of $G$: there are $\frac{\#\operatorname{Sur}(R^n,G)}{|\operatorname{Aut}G|}$ of them. For each, $|K|=\frac{p^{en}}{|G|}$, and since $p^{e-1}G=0$, $K\supseteq p^{e-1}R^n$, which is the whole $p$-torsion of $R^n$; so $K$ has $p$-torsion of dimension $n$, hence $p$-rank $n$ (for a finite abelian $p$-group, the $p$-torsion and $K/pK$ have the same dimension, the number of cyclic factors). The matrices $M$ with $MR^n=K$ are those whose $n$ columns generate $K$; by Proposition~\ref{prop:freesur} applied to $K$ (whose surjections from $R^n$ are exactly such generating $n$-tuples), there are $|K|^n\prod_{i=0}^{n-1}(1-p^{i-n})=|K|^n\prod_{j=1}^n(1-p^{-j})$ of them. Altogether
\[
\mathbb P(\operatorname{cok}M\cong G)=\frac1{p^{en^2}}\cdot\frac{|G|^n\prod_{i<r}(1-p^{i-n})}{|\operatorname{Aut}G|}\cdot\frac{p^{en^2}}{|G|^n}\prod_{j=1}^n(1-p^{-j}).\qedhere
\]
\end{proof}

\begin{corollary}[the Cohen--Lenstra law]\label{cor:clmeasure}
(a) $\sum_G\frac{c_\infty}{|\operatorname{Aut}G|}=1$, the sum over all finite abelian $p$-groups. The probability measure $\mu_{CL}(G)=\frac{c_\infty}{|\operatorname{Aut}G|}$ is the \emph{Cohen--Lenstra law}.
(b) If $X$ has law $\mu_{CL}$, then $\mathbb E\#\operatorname{Sur}(X,G)=1$ for every $G$.
(c) For $X_n=\operatorname{cok}M$ as in Theorem~\ref{thm:fw} and $G$ of exponent dividing $p^e$, $\mathbb P(\operatorname{rank}X_n\ge t)\le\frac1{|GL_t(\mathbb F_p)|}\le c_\infty^{-1}p^{-t^2}$.
\end{corollary}
\begin{proof}
(c) If $X_n$ has $p$-rank at least $t$, it has at least $|GL_t(\mathbb F_p)|$ surjections onto $\mathbb F_p^t$; so by Theorem~\ref{thm:cokmoments} and Markov's inequality, $\mathbb P(\operatorname{rank}X_n\ge t)\le\frac1{|GL_t(\mathbb F_p)|}$, and $|GL_t(\mathbb F_p)|=p^{t^2}\prod_{i\le t}(1-p^{-i})\ge c_\infty p^{t^2}$.

(a) By Fatou's lemma and Theorem~\ref{thm:fw},\[\sum_{G:\,p^{e-1}G=0}\frac{c_\infty}{|\operatorname{Aut}G|}\le1\]for every $e$, so the full sum $S$ is at most $1$. Conversely, $X_n$ has exponent $p^e$ only if it surjects onto $\Z/p^e$, which has $p^{e-1}(p-1)$ automorphisms; by Markov, this has probability at most $p^{1-e}$. By (c), all but probability $c_\infty^{-1}p^{-t^2}$ of the law of $X_n$ sits on the finitely many groups of rank $<t$ and exponent dividing $p^{e-1}$, uniformly in $n$. Letting $n\to\infty$, $S\ge1-p^{1-e}-c_\infty^{-1}p^{-t^2}$, for all $e$ and $t$.

(b) Fix $G$, let $e$ exceed its exponent by at least $1$, and $X_n$ as above. Then $\#\operatorname{Sur}(X_n,G)\le|G|^{\operatorname{rank}X_n}$, and by (c) $\mathbb E[|G|^{2\operatorname{rank}X_n}]\le\sum_t|G|^{2t}c_\infty^{-1}p^{-(t-1)^2}=:B_G$, independent of $n$ and $e$. Split $\mathbb E\#\operatorname{Sur}(X_n,G)$ according to whether $p^{e-1}X_n=0$. The part where it is not is at most $\sqrt{B_Gp^{1-e}}$ by Cauchy--Schwarz. The other part is $\sum_{H:p^{e-1}H=0}\mathbb P(X_n\cong H)\#\operatorname{Sur}(H,G)$, which converges as $n\to\infty$ to $\sum_{H:p^{e-1}H=0}\mu_{CL}(H)\#\operatorname{Sur}(H,G)$, by Theorem~\ref{thm:fw} and dominated convergence (the tail over $\operatorname{rank}H\ge t$ is at most $\sqrt{B_G}\sqrt{c_\infty^{-1}p^{-t^2}}$, uniformly in $n$). By Theorem~\ref{thm:cokmoments},\[\Bigl|\sum_{H:\,p^{e-1}H=0}\mu_{CL}(H)\#\operatorname{Sur}(H,G)-1\Bigr|\le\sqrt{B_Gp^{1-e}},\]and $e\to\infty$ gives $\mathbb E_{\mu_{CL}}\#\operatorname{Sur}(X,G)=1$ by monotone convergence.
\end{proof}

\begin{figure}[ht]
\centering\fig{ch13-cokernels}
\caption{The Cohen--Lenstra law for $p=3$ (bars, $\frac{c_\infty}{|\operatorname{Aut}G|}$) against the cokernels of $20\,000$ random $12\times12$ matrices over $\Z/81$ (dots). The trivial group has probability $c_\infty=0.5601$ and $\Z/3$ half as much; $\Z/9$, with $6$ automorphisms, is eight times as likely as $\Z/3\times\Z/3$, with $48$.}
\label{fig:cokernels}
\end{figure}

The law has striking features (Exercises~\ref{ex:clcyclic} and~\ref{ex:clrank}): the $p$-rank $r$ has probability $\frac{c_\infty p^{-r^2}}{\prod_{i=1}^r(1-p^{-i})^2}$, so large ranks are extremely rare; the group is cyclic with probability $c_\infty(1+\frac p{(p-1)^2})$, which is $0.980$ for $p=3$; and yet $\mathbb E|X|=\infty$ (Exercise~\ref{ex:clmean}), because $\Z/p^a$ has probability $\frac{c_\infty}{p^{a-1}(p-1)}$ and size $p^a$.

\section{Moments determine the law}

A real random variable is not always determined by its moments: the lognormal law is the classical counterexample, with moments $e^{k^2/2}$. The moments of random groups are of the same size. For $G=(\Z/p)^k$, Lemma~\ref{lem:homcount} gives $\#\operatorname{Hom}(X,G)=(p^{\operatorname{rank}X})^k$, so the Hom-moments of $X$ are the ordinary moments of $Y=p^{\operatorname{rank}X}$, and for the symmetric matrices of Tome~15 they grow like $p^{k^2/2}$. What saves the day is that $Y$ only takes the values $1,p,p^2,\dots$. The following lemma, due to Melanie Wood and extending a one-variable construction of Roger Heath-Brown, builds analytic functions that vanish at all but one of the possible values.

For a non-increasing vector $\mu=(\mu_1,\dots,\mu_e)$ of non-negative integers, let\[Z(\mu)=\bigl(p^{\mu_1},p^{\mu_1+\mu_2},\dots,p^{\mu_1+\dots+\mu_e}\bigr).\] We order such vectors lexicographically. For each $\nu$, the set of $\mu<\nu$ is finite, since $\mu_1\le\nu_1$ and all $\mu_i\le\mu_1$.

\begin{lemma}[Wood, 2017]\label{lem:woodfunctions}
Let $b=(b_1\ge\dots\ge b_e\ge0)$. There is an entire function $H_b(z)=\sum_{d\in\N^e}a_dz_1^{d_1}\cdots z_e^{d_e}$ on $\mathbb C^e$ and a constant $E_b$ such that: $a_d=0$ unless $d_2+\dots+d_e\le b_1$; $|a_d|\le E_bp^{-b_1d_1-\frac{d_1(d_1+1)}2}$; $H_b(Z(\mu))=0$ for every $\mu>b$; and $H_b(Z(b))\ne0$.
\end{lemma}
\begin{proof}
Let $\sigma_i=b_1+\dots+b_i$, and
\[
H_b(z)=\prod_{j>b_1}\Bigl(1-\frac{z_1}{p^j}\Bigr)\cdot\prod_{i=2}^e\ \prod_{j=\sigma_i+1}^{\sigma_{i-2}+2b_{i-1}}\Bigl(1-\frac{z_i}{p^j}\Bigr),
\]
with $\sigma_0=0$. The first factor $G(z_1)$ is an entire function, the product converging absolutely; the second is a polynomial of degree $\sigma_{i-2}+2b_{i-1}-\sigma_i=b_{i-1}-b_i$ in each $z_i$, so the total degree in $z_2,\dots,z_e$ is at most $b_1-b_e\le b_1$. Expanding $\prod_{j\ge1}(1-wq^j)$ with $q=\frac1p$ and $w=p^{-b_1}z_1$, the coefficient of $w^n$ is $\pm\sum_{j_1<\dots<j_n}q^{j_1+\dots+j_n}$; writing $j_i=i+\mu_i$ with $0\le\mu_1\le\dots\le\mu_n$, this sum is $q^{n(n+1)/2}\sum_\mu q^{\mu_1+\dots+\mu_n}\le q^{n(n+1)/2}\prod_{i=1}^n(1-q^i)^{-1}$, a non-decreasing $\mu$ being a partition into parts at most $n$. So the coefficients of $G$ are bounded by $c_\infty^{-1}p^{-b_1n-\frac{n(n+1)}2}$, and the bound on $a_d$ follows, $E_b$ being $c_\infty^{-1}$ times the largest coefficient of the polynomial part.

Let $\mu>b$, and let $t\ge0$ be the number of initial coordinates where $\mu$ and $b$ agree, so $\mu_{t+1}>b_{t+1}$. If $t=0$, $\mu_1\ge b_1+1$ and $G(p^{\mu_1})=0$. If $t\ge1$, the exponent $s=\mu_1+\dots+\mu_{t+1}$ of $z_{t+1}$ satisfies $s\ge\sigma_{t+1}+1$, and, since $\mu_{t+1}\le\mu_t=b_t$, $s\le\sigma_{t-1}+2b_t$; so a factor $1-\frac{z_{t+1}}{p^s}$ vanishes. Finally, at $Z(b)$, $z_i=p^{\sigma_i}$ while every factor involving $z_i$ has $j>\sigma_i$: none vanishes, and the infinite product $G(p^{b_1})=\prod_{k\ge1}(1-p^{-k})$ is non-zero.
\end{proof}

\begin{theorem}[Wood, 2017]\label{thm:woodunique}
Let $e\ge1$, and let $X$, $Y$ be random finite abelian $p$-groups. Suppose there is a constant $C$ such that, for every group $G$ of exponent dividing $p^e$,
\[
\mathbb E\#\operatorname{Sur}(X,G)=\mathbb E\#\operatorname{Sur}(Y,G)\le C|\wedge^2G|.
\]
Then $X/p^eX$ and $Y/p^eY$ have the same law. If this holds for every $e$, then $X$ and $Y$ have the same law.
\end{theorem}
\begin{proof}
\emph{Hom-moments.} Since every homomorphism is onto its image, $\#\operatorname{Hom}(X,G)=\sum_{H\le G}\#\operatorname{Sur}(X,H)$, so the Hom-moments of $X$ and $Y$ agree, and by Lemma~\ref{lem:subgroupsum} they are at most $C(\frac3{c_\infty})^e|\wedge^2G|$. Index the groups of exponent dividing $p^e$ by their conjugate types $\lambda=(\lambda_1\ge\dots\ge\lambda_e\ge0)$. By Lemma~\ref{lem:homcount}, with $x_\mu=\mathbb P(X/p^eX\text{ has conjugate type }\mu)$ and $F=C(\frac3{c_\infty})^e$,
\[
C_\lambda:=\sum_\mu x_\mu p^{\sum_k\lambda_k\mu_k}=\mathbb E\#\operatorname{Hom}(X,G_\lambda)\le Fp^{\sum_k\binom{\lambda_k}2},\tag{13.1}
\]
because $\operatorname{Hom}(X,G)=\operatorname{Hom}(X/p^eX,G)$ for $G$ of exponent dividing $p^e$. The same holds for $y_\mu$, the probabilities for $Y$, with the same $C_\lambda$. We show by induction along the lexicographic order that $x_\nu$ is determined by the $C_\lambda$; then $x_\nu=y_\nu$.

\emph{The lower-triangular system.} Fix $\nu$, take $H_\nu$ from Lemma~\ref{lem:woodfunctions}, and for $\lambda$ put $A_\lambda=a_d$ with $d_k=\lambda_k-\lambda_{k+1}$ ($\lambda_{e+1}=0$); this is a bijection between conjugate types and $d\in\N^e$, with $\lambda_k=d_k+\dots+d_e$. Then $\sum_k\lambda_k\mu_k=\sum_id_i(\mu_1+\dots+\mu_i)$, so for each $\mu$
\[
\sum_\lambda A_\lambda p^{\sum_k\lambda_k\mu_k}=\sum_da_d\prod_i\bigl(p^{\mu_1+\dots+\mu_i}\bigr)^{d_i}=H_\nu(Z(\mu)).
\]
\emph{Convergence.} We claim $\sum_\lambda|A_\lambda|C_\lambda<\infty$. Only $d$ with $s=d_2+\dots+d_e\le\nu_1$ contribute, finitely many choices of $(d_2,\dots,d_e)$; for each, $\lambda_1=d_1+s$ and $\lambda_2,\dots,\lambda_e\le s$. By the lemma and (13.1), the term is at most a constant times $p^{-\nu_1d_1-\frac{d_1(d_1+1)}2+\binom{d_1+s}2}=p^{d_1(s-\nu_1-1)+\binom s2}\le p^{-d_1+\binom s2}$, which is summable over $d_1$.

\emph{Conclusion.} As all $x_\mu\ge0$, the double series $\sum_\lambda\sum_\mu A_\lambda x_\mu p^{\sum_k\lambda_k\mu_k}$ converges absolutely, and summing in both orders,
\[
\sum_\lambda A_\lambda C_\lambda=\sum_\mu x_\mu H_\nu(Z(\mu))=x_\nu H_\nu(Z(\nu))+\sum_{\mu<\nu}x_\mu H_\nu(Z(\mu)),
\]
since $H_\nu(Z(\mu))=0$ for $\mu>\nu$. The sum over $\mu<\nu$ is finite and, by induction, determined by the $C_\lambda$; and $H_\nu(Z(\nu))\ne0$. So $x_\nu$ is determined.

For the last statement: for a finite group $A$, $A/p^eA=A$ as soon as $p^e$ exceeds its exponent, so $\mathbb P(X\cong A)=\lim_e\mathbb P(X/p^eX\cong A)$ (the events $\{X/p^eX\cong A\}$ decrease to $\{X\cong A\}$ when $p^e$ exceeds the exponent of $A$). The same holds for $Y$.
\end{proof}

The bound $|\wedge^2G|$ is exactly what the convergence step can absorb: the exponent $d_1(s-\nu_1-1)$ has no room to spare. It covers both cases of interest, moments equal to $1$ (cokernels of random matrices, class groups of imaginary quadratic fields) and moments equal to $|\wedge^2G|$ (cokernels of random \emph{symmetric} matrices, Jacobians of random graphs, Tome~15).

\begin{corollary}\label{cor:clcharacterize}
A random finite abelian $p$-group $X$ has the Cohen--Lenstra law if and only if $\mathbb E\#\operatorname{Sur}(X,G)=1$ for every $G$. More generally, if $X_n$ are random groups with $\mathbb E\#\operatorname{Sur}(X_n,G)\to1$ for every $G$, then $\mathbb P(X_n\cong G)\to\mu_{CL}(G)$ for every $G$.
\end{corollary}
\begin{proof}
The first statement is Theorem~\ref{thm:woodunique} with $Y$ of law $\mu_{CL}$ (Corollary~\ref{cor:clmeasure}(b)). For the second, fix $e$. As in Corollary~\ref{cor:clmeasure}(c), $\mathbb P(\operatorname{rank}X_n\ge t)\le\frac{\mathbb E\#\operatorname{Sur}(X_n,\mathbb F_p^t)}{|GL_t(\mathbb F_p)|}\le\frac2{c_\infty}p^{-t^2}$ for $n$ large, so the laws of $X_n/p^eX_n$ are tight, and the moments are uniformly integrable as in the proof of Corollary~\ref{cor:clmeasure}(b). Any subsequential limit $Z$ in law of $X_n/p^eX_n$ therefore has $\mathbb E\#\operatorname{Sur}(Z,G)=1$ for $G$ of exponent dividing $p^e$, and by Theorem~\ref{thm:woodunique} its law is that of $X/p^eX$ with $X\sim\mu_{CL}$. Every subsequence having a further subsequence with this limit, $X_n/p^eX_n$ converges in law to it, and for $G$ of exponent dividing $p^{e-1}$, $\{X_n\cong G\}$ differs from $\{X_n/p^eX_n\cong G\}$ by an event of probability at most $\frac{\mathbb E\#\operatorname{Sur}(X_n,\Z/p^e)}{p^{e-1}(p-1)}\to\frac1{p^{e-1}(p-1)}$. Let $n\to\infty$, then $e\to\infty$.
\end{proof}

\section{Class groups}

\medskip\noindent\textbf{The Cohen--Lenstra conjecture (1984).} \emph{Let $p$ be an odd prime. As $-d$ runs over the fundamental discriminants of imaginary quadratic fields ordered by $d$, the $p$-parts $\operatorname{Cl}(-d)_p$ of the class groups become distributed according to $\mu_{CL}$: for every finite abelian $p$-group $G$,
\[
\lim_{X\to\infty}\frac{\#\{d\le X:\operatorname{Cl}(-d)_p\cong G\}}{\#\{d\le X\}}=\frac{c_\infty}{|\operatorname{Aut}G|}.
\]}

By Corollary~\ref{cor:clcharacterize}, the conjecture is equivalent to: for every $G$, the average of $\#\operatorname{Sur}(\operatorname{Cl}(-d),G)$ over $d\le X$ tends to $1$. The only case proved is $G=\Z/3$: the average number of elements of order $3$ in $\operatorname{Cl}(-d)$, plus one, is $\#\operatorname{Hom}(\operatorname{Cl}(-d),\Z/3)$, whose average is $2=1+1$, a theorem of Davenport and Heilbronn (Chapter~\ref{ch:davenportheilbronn}). The prime $2$ is excluded because genus theory makes the $2$-rank of $\operatorname{Cl}(-d)$ equal to $\omega(d)-1$, which is not bounded in law; the factor $u=1$ for real quadratic fields (Cohen and Lenstra predicted $\frac{c'}{|G||\operatorname{Aut}G|}$ there) reflects the unit group.

\begin{figure}[ht]
\centering\fig{ch13-classnumbers}
\caption{The proportion of fundamental discriminants $-d$, $d\le X$, with $p\mid h(-d)$, for $p=3$, $5$, $7$ (from bottom to top: $7$, $5$, $3$), computed from $h(-d)$ for all $d\le10^7$, against the Cohen--Lenstra predictions $1-c_\infty$ (dashed): $0.4399$, $0.2397$, $0.1632$. The convergence is slow, with a secondary term of size about $X^{-1/6}$ for $p=3$, discovered by Roberts and proved by Bhargava, Shankar and Tsimerman, and by Taniguchi and Thorne.}
\label{fig:classnumbers}
\end{figure}

\begin{think}
The Cohen--Lenstra law gives the trivial group probability $c_\infty$. Compute $\prod_{p\text{ odd}}c_\infty(p)$, assuming (wrongly, but instructively) that the $p$-parts are independent for different $p$. What does it predict for the proportion of imaginary quadratic fields with odd part of the class number equal to $1$? Is that plausible, given that $h(-d)\approx\sqrt d$ on average?
\end{think}

\begin{remark}
These are quoted for culture; see the Sources. (a) \emph{Function fields.} Jordan Ellenberg, Akshay Venkatesh and Craig Westerland (2016) proved the analogue of the conjecture for quadratic extensions of $\mathbb F_q(t)$, in the limit where $q\to\infty$ suitably, by computing the moments as point counts on Hurwitz spaces and using homological stability. (b) \emph{The prime $2$.} For the $4$-ranks and higher $2$-power parts, where genus theory fixes the $2$-rank, analogous predictions (Gerth) were proved by Heath-Brown for $4$-ranks of related Selmer groups and by Fouvry and Klüners for $4$-ranks of class groups, with moments and the one-variable form of Lemma~\ref{lem:woodfunctions}; Alexander Smith proved the full distribution of the $2^\infty$-parts. (c) \emph{Universality.} Wood proved that the cokernels of random matrices whose entries are independent but not uniform (any non-degenerate distribution) have the same limit law: the Cohen--Lenstra law is a central limit theorem for groups.
\end{remark}

\begin{beauty}
The Cohen--Lenstra principle, ``weight by the inverse of the number of symmetries'', is the same principle that makes the uniform measure on labelled objects, seen up to isomorphism, give each isomorphism class a weight $\frac1{|\operatorname{Aut}|}$: it is what one gets by forgetting the labels of a random structure. A random matrix is labelled data, its cokernel is the structure seen up to isomorphism. And the moments $\mathbb E\#\operatorname{Sur}(X,G)=1$ say something even simpler: a random group knows nothing special about any given group $G$. Every target is hit once on average.
\end{beauty}

\section{Exercises}

\begin{exercise}[$\star$]\label{ex:clcyclic}
Show that under $\mu_{CL}$ the probability that $X$ is cyclic (including trivial) is $c_\infty(1+\frac p{(p-1)^2})$. Evaluate it for $p=3$, and compare with the probability $c_\infty$ that $X$ is trivial.
\end{exercise}

\begin{exercise}[$\star\star$]\label{ex:clrank}
(a) Show that the number of $n\times n$ matrices over $\mathbb F_p$ of rank $n-r$ is $\frac{\prod_{i=r}^{n-1}(p^n-p^i)^2}{\prod_{i=r}^{n-1}(p^{n-r}-p^{i-r})}$, by counting pairs (image, map). (b) Deduce that the corank of a uniform random $n\times n$ matrix over $\mathbb F_p$ converges in law to $\frac{c_\infty p^{-r^2}}{\prod_{i=1}^r(1-p^{-i})^2}$, and that this is the law of the $p$-rank under $\mu_{CL}$.
\end{exercise}

\begin{exercise}[$\star$]\label{ex:clmean}
Show that under $\mu_{CL}$, $\mathbb E|X|=\infty$, but $\mathbb E[p^{\operatorname{rank}X}]=2$ and $\mathbb E[\#X[p]]=2$, where $X[p]$ is the $p$-torsion.
\end{exercise}

\begin{exercise}[$\star\star$]\label{ex:clsur}
Verify directly, without matrices, that $\sum_{a\ge0}\mu_{CL}(\Z/p^a)\cdot\#\operatorname{Sur}(\Z/p^a,\Z/p)$ is not $1$, and explain why this does not contradict Corollary~\ref{cor:clmeasure}(b).
\end{exercise}

\begin{exercise}[$\star\star$]\label{ex:hbuniq}
(The case $e=1$, after Heath-Brown.) Let $K$ be a random non-negative integer, and suppose that $\mathbb E\bigl[\prod_{i=0}^{k-1}(p^K-p^i)\bigr]\le Cp^{\binom k2}$ for all $k$. Show, with the $q$-binomial inversion formula $\mathbf 1_{K=j}=\sum_{k\ge j}(-1)^{k-j}p^{\binom{k-j}2}\binom kj_p\binom Kk_p$, that these moments determine the law of $K$.
\end{exercise}

\section{Hints}

\begin{hint}{ex:clcyclic}
$|\operatorname{Aut}(\Z/p^a)|=p^{a-1}(p-1)$; sum a geometric series.
\end{hint}
\begin{hint}{ex:clrank}
(a) Choose the image, a subspace $V$ of dimension $k=n-r$, then a surjection $\mathbb F_p^n\to V$. (b) Divide by $p^{n^2}$, and use $\#\{V\}=\frac{\prod_{i<k}(p^n-p^i)}{\prod_{i<k}(p^k-p^i)}$.
\end{hint}
\begin{hint}{ex:clmean}
\[\#X[p]=p^{\operatorname{rank}X}=\#\operatorname{Hom}(X,\Z/p)=1+\#\operatorname{Sur}(X,\Z/p).\]
\end{hint}
\begin{hint}{ex:clsur}
Non-cyclic groups also surject onto $\Z/p$.
\end{hint}
\begin{hint}{ex:hbuniq}
$\prod_{i<k}(p^K-p^i)=|GL_k(\mathbb F_p)|\binom Kk_p$, and $\binom kj_p\le c_\infty^{-1}p^{j(k-j)}$. Check that the series of absolute values converges, then exchange the sums.
\end{hint}

\section{Solutions}

\begin{solution}{ex:clcyclic}
\idea Cyclic groups have few automorphisms, so they dominate.
\writeup $\mu_{CL}(\text{cyclic})=c_\infty\bigl(1+\sum_{a\ge1}\frac1{p^{a-1}(p-1)}\bigr)=c_\infty\bigl(1+\frac1{p-1}\cdot\frac1{1-1/p}\bigr)=c_\infty\bigl(1+\frac p{(p-1)^2}\bigr)$. For $p=3$: $0.5601\times1.75=0.9802$. So $98\%$ of the $3$-parts should be cyclic, although $44\%$ are non-trivial.
\end{solution}

\begin{solution}{ex:clrank}
\idea A matrix is its image together with a surjection onto it.
\writeup (a) A matrix of rank $k$ is a surjective linear map from $\mathbb F_p^n$ onto its image $V$, of dimension $k$; for a given $V$ there are $\prod_{i<k}(p^n-p^i)$ surjections (Proposition~\ref{prop:freesur} with $e=1$), and the number of $V$ is $\frac{\prod_{i<k}(p^n-p^i)}{\prod_{i<k}(p^k-p^i)}$ (ordered bases of $V$ in $\mathbb F_p^n$, divided by ordered bases of $\mathbb F_p^k$). The product is $\frac{\prod_{i<k}(p^n-p^i)^2}{\prod_{i<k}(p^k-p^i)}$, which is the formula with $k=n-r$, after reindexing. (b) Dividing by $p^{n^2}$: $\prod_{i<k}(p^n-p^i)^2=p^{2nk}\prod_{i<k}(1-p^{i-n})^2$ and $\prod_{i<k}(p^k-p^i)=p^{k^2}\prod_{j=1}^k(1-p^{-j})$. So the probability is $p^{2nk-k^2-n^2}\frac{\prod_{j=r+1}^n(1-p^{-j})^2}{\prod_{j=1}^k(1-p^{-j})}=p^{-r^2}\frac{\prod_{j=r+1}^n(1-p^{-j})^2}{\prod_{j=1}^{n-r}(1-p^{-j})}$, which tends to $p^{-r^2}\frac{c_\infty^2/\prod_{j\le r}(1-p^{-j})^2}{c_\infty}$. The $p$-rank of $\operatorname{cok}M$ is the corank of $M\bmod p$, since $\operatorname{cok}M/p\operatorname{cok}M=\operatorname{cok}(M\bmod p)$; by Theorem~\ref{thm:fw} and Corollary~\ref{cor:clmeasure}, the $p$-rank of a large random cokernel tends in law to that under $\mu_{CL}$, so the two laws agree.
\end{solution}

\begin{solution}{ex:clmean}
\idea Moments of size $1$ for every $G$ do not prevent a heavy tail for $|X|$.
\writeup $\mathbb E|X|\ge\sum_{a\ge1}p^a\mu_{CL}(\Z/p^a)=\sum_a\frac{c_\infty p}{p-1}=\infty$. On the other hand $\#X[p]=p^{\operatorname{rank}X}=\#\operatorname{Hom}(X,\Z/p)=1+\#\operatorname{Sur}(X,\Z/p)$, whose expectation is $1+1=2$ by Corollary~\ref{cor:clmeasure}(b).
\end{solution}

\begin{solution}{ex:clsur}
\idea The moment sees all groups, not only the cyclic ones.
\writeup For $a\ge1$, $\#\operatorname{Sur}(\Z/p^a,\Z/p)=p-1$, so the sum is $(p-1)\sum_{a\ge1}\frac{c_\infty}{p^{a-1}(p-1)}=\frac{c_\infty p}{p-1}$, which is $0.840$ for $p=3$, not $1$. The missing $0.160$ comes from the non-cyclic groups: $\Z/p\times\Z/p$ alone has $p^2-1$ surjections onto $\Z/p$ and probability $\frac{c_\infty}{p(p-1)^2(p+1)}$, contributing $\frac{c_\infty}{p(p-1)}=0.093$ for $p=3$, and so on.
\end{solution}

\begin{solution}{ex:hbuniq}
\idea Gaussian binomials of $p^K$ are to $p$-ranks what binomial coefficients are to integers.
\writeup Write $m_k=\mathbb E\binom Kk_p=\frac{\mathbb E\prod_{i<k}(p^K-p^i)}{|GL_k(\mathbb F_p)|}\le\frac{Cp^{\binom k2}}{c_\infty p^{k^2}}=C'p^{-\binom{k+1}2}$. By the inversion formula, $\mathbb P(K=j)=\mathbb E\sum_{k\ge j}(-1)^{k-j}p^{\binom{k-j}2}\binom kj_p\binom Kk_p$, and the exchange of $\mathbb E$ and $\sum$ is justified if $\sum_kp^{\binom{k-j}2}\binom kj_pm_k<\infty$. With $k=j+t$, the term is at most $c_\infty^{-1}C'p^{\binom t2+jt-\binom{j+t+1}2}=c_\infty^{-1}C'p^{-t-\binom{j+1}2}$, summable. So $\mathbb P(K=j)=\sum_k(-1)^{k-j}p^{\binom{k-j}2}\binom kj_pm_k$ is determined by the moments. (The inversion formula itself follows from the $q$-binomial theorem $\sum_{i}(-1)^ip^{\binom i2}\binom ni_p=\mathbf 1_{n=0}$, applied with $n=K-j$ after writing $\binom kj_p\binom Kk_p=\binom Kj_p\binom{K-j}{k-j}_p$.) This is the case $e=1$ of Theorem~\ref{thm:woodunique}, where the lexicographic induction is not needed.
\end{solution}

\section*{Sources and further reading}

\begin{itemize}
\item \emph{The heuristic.} H.~Cohen and H.~W.~Lenstra, Jr., ``Heuristics on class groups of number fields'', in \emph{Number Theory, Noordwijkerhout 1983}, Lecture Notes in Math.\ 1068, Springer, 1984, 33--62.
\item \emph{Random matrices.} E.~Friedman and L.~C.~Washington, ``On the distribution of divisor class groups of curves over a finite field'', in \emph{Théorie des Nombres (Quebec, 1987)}, de Gruyter, 1989, 227--239.
\item \emph{Moments.} M.~M.~Wood, ``The distribution of sandpile groups of random graphs'', \emph{J.\ Amer.\ Math.\ Soc.}\ 30 (2017), 915--958, Section~8, whose Lemma~8.1 and Theorem~8.2 we follow; M.~M.~Wood, ``Random integral matrices and the Cohen--Lenstra heuristics'', \emph{Amer.\ J.\ Math.}\ 141 (2019), 383--398; M.~M.~Wood, ``Probability theory for random groups arising in number theory'', \emph{Proceedings of the ICM 2022}, EMS Press, 2023, a survey of the whole subject; D.~R.~Heath-Brown, ``The size of Selmer groups for the congruent number problem, II'', \emph{Invent.\ Math.}\ 118 (1994), 331--370, Lemma~17.
\item \emph{Proved cases.} J.~S.~Ellenberg, A.~Venkatesh and C.~Westerland, ``Homological stability for Hurwitz spaces and the Cohen--Lenstra conjecture over function fields'', \emph{Ann.\ of Math.}\ 183 (2016), 729--786; É.~Fouvry and J.~Klüners, ``On the 4-rank of class groups of quadratic number fields'', \emph{Invent.\ Math.}\ 167 (2007), 455--513; A.~Smith, ``$2^\infty$-Selmer groups, $2^\infty$-class groups, and Goldfeld's conjecture'', preprint, arXiv:1702.02325.
\item \emph{Secondary terms.} M.~Bhargava, A.~Shankar and J.~Tsimerman, ``On the Davenport--Heilbronn theorems and second order terms'', \emph{Invent.\ Math.}\ 193 (2013), 439--499; T.~Taniguchi and F.~Thorne, ``Secondary terms in counting functions for cubic fields'', \emph{Duke Math.\ J.}\ 162 (2013), 2451--2508.
\end{itemize}
